\src{178}
\[
\symup{\Extsh_{\Osh_X}^{p}(F,L)
=L\otimes(\underline{\Omega}_X^{n})'\otimes
  \Extsh_{\Osh_X}^{p}(F,\underline{\Omega}_X^{n})
=L\otimes(\underline{\Omega}_X^{n})'\otimes
  \Homsh_{\Osh_X}(F,\underline{\omega}_Y^{q})
=\Homsh_{\Osh_X}\!\left(
  F\otimes L'\otimes\underline{\Omega}_X^{n},\underline{\omega}_Y^{q}\right)}
\]
où $\symup{q=n-p}$, d'où la formule (4.1 bis).

Faisons en particulier $\symup{F=\Osh_Y}$, $\symup{L=\underline{\Omega}_X^{p}}$, on trouve, compte tenu
que $\symup{\underline{\Omega}_X^{n}\otimes(\underline{\Omega}_X^{p})'=\underline{\Omega}_X^{n-p}}$, un isomorphisme
canonique
\begin{equation}\tag{4.2}
  \symup{\operatorname{Ext}_{\Osh_X}^{p}(X;\Osh_Y,\underline{\Omega}_X^{p})
  =\operatorname{Hom}_{\Osh_X}
  (X;\underline{\Omega}_X^{n-p},\underline{\omega}_Y^{n-p})}
\end{equation}
Supposons maintenant $\symup{Y}$ \uline{non singulière} pour simplifier, donc
$\symup{\underline{\omega}_Y^{n-p}=\underline{\Omega}_Y^{n-p}}$. On a un homomorphisme naturel de
$\symup{\underline{\Omega}_X^{n-p}}$ dans $\symup{\underline{\Omega}_Y^{n-p}}$, d'où une section canonique $\symup{s_Y}$ du
faisceau $\symup{\Extsh_{\Osh_X}^{p}(\Osh_Y,\underline{\Omega}_X^{p})}$, que nous appellerons, si
toutes les composantes de $\symup{Y}$ sont de dimension $\symup{n-p}$, \uline{section
fondamentale} du faisceau $\symup{\Extsh_{\Osh_X}^{p}(\Osh_Y,\underline{\Omega}_X^{p})}$. Cette
section définit en vertu de (4.1) un élément de
$\symup{\operatorname{Ext}_{\Osh_X}^{p}(X;\Osh_Y,\underline{\Omega}_X^{p})}$. Or l'homomorphisme
naturel $\symup{\Osh_X\longrightarrow\Osh_Y}$ définit un homomorphisme
\[
  \symup{\operatorname{Ext}_{\Osh_X}^{p}(X;\Osh_Y,\underline{\Omega}_X^{p})
  \longrightarrow
  \operatorname{Ext}_{\Osh_X}^{p}(X;\Osh_X,\underline{\Omega}_X^{p})
  =H^{p}(X,\underline{\Omega}_X^{p}).}
\]
On obtient ainsi un élément de $\symup{H^{p}(X,\underline{\Omega}_X^{p})}$, noté $\symup{P_X(Y)}$ et appelé
\uline{classe de cohomologie de} $\symup{Y}$ \uline{dans} $\symup{X}$. Elle est donc déduite de la section
$\symup{s_Y}$ de $\symup{\Extsh_{\Osh_X}^{p}(\Osh_Y,\underline{\Omega}_X^{p})}$ par le diagramme
d'homomorphismes suivants :
\begin{equation}\tag{4.3}
\symup{H^{p}(X,\underline{\Omega}_X^{p})
=\operatorname{Ext}_{\Osh_X}^{p}(X;\Osh_X,\underline{\Omega}_X^{p})
\longleftarrow\operatorname{Ext}_{\Osh_X}^{p}(X;\Osh_Y,\underline{\Omega}_X^{p})
\xrightarrow{\ \sim\ }
H^{0}\!\left(X,\Extsh_{\Osh_X}^{p}(\Osh_Y,\underline{\Omega}_X^{p})\right)}
\end{equation}

Appelons \uline{cycle non singulier} de dimension $\symup{n-p}$ un élément du groupe
abélien libre engendré par les sous-variétés irréductibles non singulières de
dimension $\symup{n-p}$ dans $\symup{X}$. Alors la fonction $\symup{Y\longrightarrow P(Y)}$ se prolonge en
un homomorphisme du groupe des cycles non singuliers de dimension $\symup{n-p}$ sur $\symup{X}$,
dans le groupe $\symup{H^{p}(X,\underline{\Omega}_X^{p})}$.

Soient $\symup{Z^{n-p}}$ et $\symup{Z'^{n-p'}}$ des cycles non singuliers de dimension $\symup{n-p}$ et
$\symup{n-p'}$, on dit qu'ils se \uline{coupent transversalement}, si toute composante
de $\symup{Z}$ coupe transversalement toute composante de $\symup{Z'}$. Alors le cycle $\symup{Z.Z'}$ est
défini, c'est un cycle non singulier de dimension $\symup{n-p-p'}$. Ceci dit, on a le

